% Figura redibujada a partir de "Imagenes/Regiones de densidad de portadores.PNG" (Tema 1).
% Ocupa lo mismo que el PNG en el documento: 201.1 pt x 144.8 pt. La posición
% de los ejes se da en píxeles del PNG original (835 x 601 px).
% Curvas calculadas para Si con N_d = 5e14 cm^-3 (donadores a 0.045 eV de E_C),
% N_C = 2.8e19 cm^-3, N_V = 1.04e19 cm^-3 a 300 K y E_g = 1.12 eV.
\documentclass[12pt,border=0pt]{standalone}
\usepackage[spanish,es-tabla]{babel}
\usepackage[utf8]{inputenc}
\usepackage{amsmath}
\usepackage{xcolor}
\usepackage{tikz}
\usepackage{pgfplots}
\usepackage{circuitikz}
\usetikzlibrary{babel, circuits.ee}
\pgfplotsset{compat=1.18}

\definecolor{applegreen}{rgb}{0.55, 0.71, 0.0}
\definecolor{airforceblue}{rgb}{0.36, 0.54, 0.66}
\definecolor{amber}{rgb}{1.0, 0.49, 0.0}

\newcommand{\diff}[3][1]{%
	\IfBlankTF{#1}{\frac{\mathrm{d}{#2}}{\mathrm{d}{#3}}}%
	{\frac{\mathrm{d}^{#1} {#2}}{\mathrm{d}{#3}^{#1}}}%
}%
\newcommand{\pdiff}[3][1]{%
	\IfBlankTF{#1}{\frac{\partial {#2}}{\partial {#3}}}%
	{\frac{\partial^{#1} {#2}}{\partial {#3}^{#1}}}%
}%

\begin{document}
\begin{tikzpicture}[x=0.24084pt, y=-0.24084pt, font=\scriptsize]
    \useasboundingbox (0,0) rectangle (835,601);
    \begin{scope}[x=1pt, y=1pt]
    % Esquina inferior izquierda de los ejes en (133, 512) px
    \begin{axis}[
        at={(32.03pt,-123.31pt)}, anchor=south west, scale only axis,
        width=163.06pt, height=115.36pt,
        xmin=0, xmax=711, ymode=log, ymin=1.622e11, ymax=1e16,
        axis x line*=bottom, axis y line*=left,
        xtick={100,200,...,700}, ytick={1e12,1e13,1e14,1e15,1e16},
        extra y ticks={5e14}, extra y tick labels={$5$},
        tick align=outside, tick style={line width=0.4pt, black},
        major tick length=2.5pt, yminorticks=false,
        every axis plot/.append style={samples=200},
        declare function={
            kT(\T) = 8.617e-5*\T;
            Nc(\T) = 2.8e19*(\T/300)^1.5;
            ni(\T) = 1.7065e19*(\T/300)^1.5*exp(-1.12/(2*kT(\T)));
            Ndmas(\T) = 1e15/(1 + sqrt(1 + 4e15*exp(0.045/kT(\T))/Nc(\T)));
            n0(\T) = Ndmas(\T)/2 + sqrt((Ndmas(\T)/2)^2 + ni(\T)^2);
        },
    ]
        \addplot[blue, ultra thick, domain=27.5:700] {n0(x)};
        \addplot[black, thick, dashed, domain=388:705] {ni(x)};
    \end{axis}
    \end{scope}
    % Corte del eje vertical
    \fill[white] (126,455) rectangle (140,477);
    \draw (125,462) -- (141,457) (125,472) -- (141,467);
    % Etiquetas
    \node[anchor=west] at (140,22) {$n_0$ (cm$^{-3}$)};
    \node at (470,586) {$T$ (K)};
    \begin{scope}[text=blue]
        \node[align=left] at (292,270) {Ionización\\parcial};
        \node at (437,137) {Extrínseca};
        \node[anchor=east] at (765,48) {Intrínseca};
    \end{scope}
    \node at (558,365) {$n_\mathrm{i}$};
\end{tikzpicture}
\end{document}
